%*****************************************
\chapter{Other functional testing methods}\label{app:comparison-test}
%*****************************************

We provide a brief description of the comparison methods used in the experiments in Chapter~\ref{ch:functional-testing}. 

\section{Frequentist methods}

The frequentist comparison methods were described in Chapter~4 of the monograph by \citet{zhang2013analysis}. They revolve around the pivotal test function \(\Delta_n(s) = \sqrt{n}\bar X_n(s)\), which is small when the null hypothesis is true. In what follows, we fix a significance level \(\alpha\in (0,1)\) and the discretized grid \(\mathcal S=\{s_1,\dots,s_G\}\).

\subsubsection*{Pointwise \(t\)-test}

The first naive idea is to construct a pointwise \(t\)-test over all points on the grid, and then reject \(H_0\) if any pointwise test rejects the local null hypothesis \(H_{0s}: \mu(s)=0\). The test statistics are
\[
z_n(s) = \frac{\Delta_n(s)}{\sqrt{\hat C_n(s,s)}} \sim t_{n-1},\quad s\in[0,1].
\]

We use the Bonferroni correction on the \(p\)-values~\citep{shaffer1995multiple} to produce a global test that accounts for the multiple comparisons problem and controls the probability of the so-called family-wise error rate, at the expense of decreased power. In particular, we reject \(H_0\) whenever the \(p\)-value of at least one of the local tests is less than \(\alpha/G\). While there are more sophisticated approaches to correct \(p\)-values \citep[e.g.][]{cox2008pointwise}, we keep the simple Bonferroni method as a baseline.

\subsubsection*{\(L^2\)-norm based tests}

We perform a global test based on the squared \(L^2\)-norm of the pivotal test function \(\Delta_n\). The proposed test statistic is
\begin{equation}\label{eq:test-statistic-l2}
T_n = \|\Delta_n\|^2_2 = n\int_0^1 \bar{X}_n^2(s)\, ds.
\end{equation}
Under \(H_0\) and Gaussian data, \(T_n\) follows a \(\chi^2\)-type mixture distribution. Using a Box-type approximation \citep{box1954some}, we have
\[
T_n \sim \gamma \chi^2_d\,  \text{ approximately}, \quad \text{where } \gamma = \frac{\tr(K^{\otimes 2})}{\tr(K)}\, \text{ and }\, d=\frac{\tr^2(K)}{\tr(K^{\otimes 2})},
\]
with \(\tr(K^{\otimes 2}) = \int_0^1 \int_0^1 K(t,s)^2 \, dt\, ds\). Since \(K\) is generally unknown, a common approach is to use the empirical covariance estimator \(\hat C_n\) to approximate \(\gamma\) and \(d\), which leads to consistent estimators \(\hat\gamma\) and \(\hat d\) \citep[Theorem~4.17]{zhang2013analysis}. The corresponding \(p\)-value for the test is given by \(\pp(\chi^2_{\hat d} \geq T_n/\hat \gamma)\). In practice, we need to discretize the functional elements to compute the test statistic and \(p\)-value:
\[
T_n = n\sum_{g=1}^G \bar{X}_n^2(s_g), \qquad \hat \gamma \approx \frac{\tr(\hat{\mathbf C}_n^2)}{\tr(\hat{\mathbf C}_n)}, \qquad \hat d \approx \frac{\tr^2(\hat{\mathbf C}_n)}{\tr(\hat{\mathbf C}_n^2)}.
\]
Note that the constant discretization step \(1/G\) that would appear from approximating the integrals cancels out in the computations, so it can be omitted for simplicity.

\begin{remark}
  If we use the improved estimator \(\hat{\mathbf K}_n^{\text{OAS}}\) to approximate \(\hat \gamma\) and \(\hat d\), the argument goes through unaffected: by Proposition~\ref{prop:sample-cov-matrix-convergence} we have \(\|\hat{\mathbf K}^{\text{OAS}}_n - \mathbf{K}\| \to 0\) a.s., and hence the traces converge almost surely to the true value.
\end{remark}

We can also consider a projection-based version of the \(L^2\) test using \acs{FPCA}, which was explored for example by \citet{horvath2012inference} in a similar testing problem. Let \(\lambda_1\geq \lambda_2\geq \cdots \geq \lambda_p>0\) be the \(p\) leading eigenvalues of the true covariance operator, with \(\phi_1,\dots,\phi_p\) the corresponding eigenfunctions (up to a sign flip). Then, we can define a test statistic analogous to~\eqref{eq:test-statistic-l2} but with the mean function projected on the space spanned by \(\phi_1,\dots,\phi_p\). Write the \acs{FPCA} scores of the sample mean as
\begin{equation}\label{eq:l2-fpca-scores}
\bar{\xi}_j = \int_0^1 \bar X_n(s)\phi_j(s)\, ds,\quad j=1,\dots,p,
\end{equation}
and consider the statistic
\begin{equation}\label{eq:statistic-l2-fpca}
T_{n,p} = n \sum_{j=1}^p \frac{\bar{\xi}_j^2}{\lambda_j},
\end{equation}
where normalization by the eigenvalues is introduced to simplify the expression of the null distribution. Indeed, by standard \acs{FPCA} theory, under \(H_0\) and the Gaussian assumption we have \(T_{n,p}\sim \chi^2_p\). Of course, we rarely have access to the true \(K\), so we construct estimates \(\{\hat\lambda_j\}\) and \(\{\hat\phi_j\}\) for the eigenvalues and eigenfunctions via the (truncated) spectral decomposition of the empirical covariance \(\hat C_n\), and then compute an approximate \(\hat T_{n,p}\) plugging in the estimated values into~\eqref{eq:l2-fpca-scores} and~\eqref{eq:statistic-l2-fpca}. Under mild assumptions (finite fourth moments, distinctness of the true eigenvalues), it can be shown that \(\hat T_{n,p}\stackrel{d}{\to}\chi^2_p\) as \(n\to\infty\) \citep[Theorem~5.3]{horvath2012inference}, so asymptotically the \(p\)-value for the test is \(\pp(\chi^2_p \geq \hat T_{n,p})\). As we said before, in practice we compute all quantities involved in the discretized grid \(\mathcal S\):
\[
\hat \lambda_j \approx \hat l_j,\quad \hat\phi_j \approx \hat{\mathbf v}_j=(\hat v_{j1},\dots, \hat v_{jG})', \quad \hat T_{n,p} \approx n\sum_{j=1}^p \frac{\left(\sum_{g=1}^G \bar{X}_n(s_g)\hat v_{jg}\right)^2}{\hat l_j}, 
\]
where \(\{\hat l_j, \hat{\mathbf v}_j\}\) are the eigenvalues and (normalized) eigenvectors of the \(G\times G\) empirical covariance matrix \(\hat{\mathbf C}_n\). We could also use the regularized version \(\hat{\mathbf K}_n^{\text{OAS}}\) to compute the eigendecomposition, and even though we did not explore the theoretical guarantees in this case, it is likely that the eigenvalues and eigenvectors of the \acs{OAS} estimator constitute sufficiently good approximations for our purposes.

The only missing aspect is the question of how many components to retain. We follow the simple approach proposed by \citet{ramsay2005functional}, who suggest selecting \(p\) so that the resulting functional principal components cumulatively explain a fixed large percentage of the variance (e.g.\ 90\%).

\subsubsection*{\(F\)-type test}

The previous \(L^2\)-based tests do not take into account the variation of the covariance function. Under \(H_0\) and Gaussian data, \(\E\|\Delta_n\|^2_2 = \tr(K)\) and \(\E[\tr(\hat{C}_n)] = \tr(K)\), so we can consider the statistic
\[
F_n = \frac{\|\Delta_n\|^2_2}{\tr(\hat C_n)} \approx \frac{n\sum_{g=1}^G \bar{X}_n^2(s_g)}{\tr(\hat{\mathbf C}_n)},
\]
whose null distribution is shown to be an \(F\)-type mixture. It can again be approximated by a two-cumulant method \citep[\S4.4]{zhang2013analysis}, so that
\[
F_n \sim F_{\hat d, (n-1)\hat d}\  \ \text{approximately}.
\]
The approximate \(p\)-value for the test is then \(\pp(F_{\hat d, (n-1)\hat d} \geq F_n)\). This method works for the \acs{OAS} estimator as well, since \(\tr(\hat{\mathbf K}_n^{\text{OAS}})=\tr(\hat{\mathbf K}_n)\).

\subsubsection*{Bootstrap procedures}

Both the \(L^2\)-based tests and the \(F\)-type test rely on asymptotic approximations, so their performance can deteriorate when the sample size \(n\) is small. For such cases we can resort to a nonparametric bootstrap scheme to approximate the null distribution of the statistics. In particular, we consider the following simple algorithm:
\begin{enumerate}
  \item Draw a bootstrap sample \(X_1^*,\dots, X_n^*\) from \(X_1,\dots, X_n\) (with replacement).
  \item Compute the quantities \(\bar{X}_n^*\) and \(\hat{\mathbf C}_n^*\) from the bootstrap sample.
  \item Construct the bootstrap statistics \(T_n^* = \|\sqrt{n}(\bar{X}_n^*(s) - \bar{X}(s))\|^2_2\) and \(F_n^* = T_n^*/\tr(\hat{\mathbf C}_n^*)\).
  \item Repeat the process many times to obtain \(R\) bootstrap statistics, and compute the bootstrap \(p\)-value as the proportion of bootstrap statistics that exceed the original statistic.
\end{enumerate}
In particular, the bootstrap \(p\)-values are
\[
p^*(T_n) = \frac{1 + \#\{T_n^* \geq T_n\}}{R+1},\qquad p^*(F_n) = \frac{1 + \#\{F_n^* \geq F_n\}}{R+1},
\]
where the original statistic is added as a positive instance to avoid the case of \(p^*=0\) \citep[see][p.~141]{davison1997bootstrap}. Note that the same procedure can be applied to the \(L^2\) test with \acs{FPCA}, recomputing the eigendecomposition of the bootstrap empirical covariance matrix \(\hat{\mathbf C}_n^*\) each time. In our implementation, we activate this bootstrap procedure with \(R=5000\) whenever \(n\leq 50\).

\section{Bayesian methods}

Consider the discretized vector \(\mathbf X_i=(X_i(s_1), \dots, X_i(s_G))'\) for \(i=1,\dots, n\), and place a multivariate normal model on these finite-dimensional observations:
\[
\mathbf X_i\mid \symbfit{\mu}, \mathbf K \sim \mathcal N_G(\symbfit \mu, \mathbf K),\quad  i=1,\dots,n,
\]
where \(\symbfit\mu=(\mu(s_1), \dots, \mu(s_G))'\) and \(\mathbf K\) is the \(G\times G\) discretized covariance on the grid. To test \(H_0: \symbfit\mu = 0\) against \(H_1: \symbfit\mu \neq 0\), we place a \(g\)-prior on the mean under \(H_1\), i.e.,
\[
\symbfit \mu \mid g, \mathbf K \sim \mathcal{N}_G(0, g \mathbf K),
\]
where \(g>0\) is a scalar hyperparameter. This prior, introduced by \citet{zellner1986assessing}, encodes the belief that the mean function varies along the same directions as the data covariance, scaled by the factor \(g\). Under this conjugate model, the Bayes factor in favor of the alternative is available in closed form after integrating out \(\symbfit \mu\) \citep[for a similar computation, see][]{liang2008mixtures}:
\begin{equation}\label{eq:log-bf10}
\log \bff_{10} = -\frac{G}{2}\log(1+ng) + \frac{n^2 g}{2(1+ng)}Q_n,
\end{equation}
where \(Q_n = \bar{\mathbf X}_n' \mathbf K^{-1}\bar{\mathbf X}_n\), \(\bar{\mathbf X}_n = \frac{1}{n}\sum_{i=1}^n \mathbf X_i\) is the discretized sample mean, and we have supposed for simplicity that \(\pi_0=1-\pi_0=1/2\).

In practice, we compute the Bayes factor using the plug-in estimators \(\hat{\mathbf C}_n\) or \(\hat {\mathbf K}_n^{\text{OAS}}\) of the covariance matrix. However, applying~\eqref{eq:log-bf10} directly on the full \(G\)-dimensional discretized data presents two noticeable difficulties:
\begin{enumerate}
  \item When \(n<G\) (the typical \acs{FDA} setting), the empirical covariance matrix \(\hat{\mathbf C}_n\) is rank-deficient, so \(\hat{\mathbf C}_n^{-1}\) and the corresponding plug-in quadratic form \(\hat Q_n\) do not exist. Even with regularization (e.g.\ by \acs{OAS} or adding a small noise to the diagonal), the resulting \(\hat Q_n\) can be numerically unstable and a poor estimate.
  \item The Bayes factor contains a penalty term that grows linearly with \(G\). This penalty can become large and drown out the signal in \(Q_n\), driving \(\bff_{10}\) towards \(H_0\) regardless of the data. This phenomenon is known as Lindley's paradox \citep{lindley1957statistical}.
\end{enumerate}

From a functional perspective, both problems are organically resolved by projecting onto the top \(p\ll G\) functional principal components, retaining the signal-carrying directions and discarding the noisy ones. As before, consider the eigendecomposition
\[
\mathbf K = \mathbf V \Sigma \mathbf V', \qquad \mathbf V = [\mathbf v_1 \mid \cdots \mid \mathbf v_p], \qquad \Sigma = \diag(l_1,\dots, l_p),
\]
where \(l_1 \geq \cdots \geq l_p > 0\) are the leading eigenvalues of \(\mathbf K\) and \(\mathbf v_1,\dots, \mathbf v_p\) are the corresponding eigenvectors. We can transform the data onto the \(p\)-dimensional score space, defining the \acs{FPCA} scores as
\[
\xi_i = \mathbf V' \mathbf X_i, \quad i=1,\dots,n.
\]
In score space, the null hypothesis can be expressed as \(H_0: \tilde{\mu}=0\), where \(\tilde{\mu}=(\tilde \mu_1,\dots, \tilde \mu_p)'\) and \(\tilde \mu_j = \symbfit \mu' \mathbf v_j\) is the projection of the true mean vector onto the \(j\)-th eigenvector. In this way, the model becomes
\[
\xi_i\mid \tilde{\mu}, \Sigma \sim \mathcal N_p(\tilde{\mu}, \Sigma), \quad i=1,\dots,n,
\]
and the corresponding \(g\)-prior under \(H_1\) is
\[
\tilde{\mu}\mid g, \Sigma \sim \mathcal N_p(0, g\Sigma).
\]
We can derive the Bayes factor for this reduced model in the same way as we did above for the fully discretized model (in fact now it is easier because the prior factorizes in independent components), obtaining
\begin{equation}\label{eq:log-bf10-fpca}
\log \bff_{10}^{\text{FPCA}} = -\frac{p}{2}\log(1+ng) + \frac{n^2 g}{2(1+ng)} Q_{n,p},
\end{equation}
where 
\[
Q_{n,p} = \sum_{j=1}^p \frac{(\bar {\mathbf X}_n' \mathbf v_j)^2}{l_j}.
\]
Observe that \(Q_{n,p}\) is essentially a discretized and rescaled version of the statistic \(T_{n,p}\) in the \(L^2\)-based \acs{FPCA} test, defined in~\eqref{eq:statistic-l2-fpca}. Naturally, in the actual implementation we approximate the eigenvalues and eigenvectors via the eigendecomposition of the empirical covariance matrix \(\hat {\mathbf C}_n\) or its regularized version, obtaining empirical estimates \(\{\hat l_j, \hat{\mathbf v}_j\}\) that we can substitute into the formulas above.

We still need to choose a suitable value for the hyperparameter \(g\). It controls the prior spread under \(H_1\): a small \(g\) concentrates the prior near the null, while a large \(g\) makes it diffuse. A fixed choice (e.g.\ \(g=1\)) is common, but data-driven selection typically performs better. As mentioned by~\citet{liang2008mixtures}, we can compute an empirical Bayes estimate of \(g\) by maximizing the marginal likelihood under \(H_1\). In particular, consider
\[
g^*_n = \operatorname*{argmax}_g\, \bff_{10}^{\text{FPCA}}(g) = \frac{Q_{n,p}}{p} - \frac{1}{n},
\]
which is obtained by differentiating~\eqref{eq:log-bf10-fpca} with respect to \(g\) and setting the derivative to zero. Then, the estimated value for \(g\) is \(\hat g_n = \max(\delta, g^*_n)\), where \(\delta>0\) is a small constant to avoid degeneracy at the boundary \(\{g=0\}\) (in the experiments we used \(\delta=10^{-3}\)).

Note that this \acs{FPCA}-based Bayesian method is simpler than the full \acs{RJMCMC} procedure described in Section~\ref{sec:rkhs-test}, and provides an exact Bayes factor without the need to approximate the posterior. However, it is not as flexible as the \acs{RKHS}-based approach, and it seems to perform a bit worse in practice.

%*****************************************
%*****************************************
%*****************************************
%*****************************************
%*****************************************
